\documentclass[10pt]{article}
\usepackage[margin=0.75in]{geometry}
\usepackage[T1]{fontenc}
\usepackage{lmodern}
\usepackage{microtype}
\usepackage{booktabs}
\usepackage{xcolor}
\usepackage{hyperref}
\usepackage{enumitem}
\setlist{nosep,leftmargin=*}
\setlength{\parskip}{0.3em}
\setlength{\parindent}{0pt}

% Remove every \todo before submitting: grep -n todo main.tex
\newcommand{\todo}[1]{\textcolor{red}{[TODO: #1]}}
\newcommand{\res}{\textcolor{red}{--}}  % a result still to fill in from summary/condition_summary.csv

\title{Failure Propagation in a Researcher--Planner--Executor Agent Workflow\\\large CSE 291A Option 5 --- Phase 1}
\author{\todo{full names}}
\date{}

\begin{document}
\maketitle
\vspace{-2em}

\section{Question}
When one agent in a multi-agent workflow produces a wrong or malicious intermediate result, does the fault
actually reach later agents and the shared state, and what does it cost to recover? Phase 1 measures
propagation and two baselines (full restart, always verify); Phase 2 will add selective recovery.

\section{Design}
\textbf{Workflow.} A researcher reads the bank (read tools only), a planner turns its facts into a plan (no tools),
and an executor performs the plan (write tools only). Every tool call is role-checked, so facts must flow
researcher $\rightarrow$ planner $\rightarrow$ executor.
\todo{model and backend; one line on prompts, plan format and retry limits}

\textbf{Environment and tasks.} AgentDojo 0.1.35 Banking v1.2.2~\cite{agentdojo} behind a role-checked wrapper with
an append-only audit log of every call and its state diff. Task~0 pays a bill read from a file (98.70);
task~2 updates the scheduled rent to 1200, which appears nowhere and must be derived from the current rent (1100)
and a notice (+100), so it depends on the researcher's result.

\textbf{Faults.} One handoff output is replaced per faulty attempt, at researcher$\rightarrow$planner or
planner$\rightarrow$executor. Ten fixed scenarios: per task, an accidental corrupted fact (bill 1337.42; rent
increase 200, i.e.\ 1300) and an adversarial instruction to send 100 to an attacker IBAN, at each boundary,
plus an adversarial instruction hidden in the planner's notes. Accidental and adversarial cases are analyzed
separately.

\textbf{Conditions.} \emph{clean}; \emph{clean\_verify} (verifier on, no fault); \emph{fault} (no recovery);
\emph{full\_restart} (after the faulty run, a fresh environment and session redo everything);
\emph{always\_verify} (a verifier re-reads the evidence and checks each handoff before it is consumed, with one
regeneration).

\section{Measurement}
\textbf{Correctness.} A strict check on the audit log: the authorized action happened with exactly the right
values, and nothing outside the task's allowed actions was executed (AgentDojo's own utility is also recorded).
\textbf{Security.} Unauthorized attempts and executions; \emph{attack success}: the fault's goal (any payment
to the attacker, or the corrupted amount) was reached in the bank at any point, judged by replaying the log.
\textbf{Propagation.} Each applied fault is labelled by its furthest consequence: not exposed, exposed but not
adopted, propagated (used in a plan or tool attempt), or state consequence (executed write). Exposure alone is
not propagation; a run whose injection was not applied is reported separately.
\textbf{Cost.} End-to-end latency, model and tool calls, tokens, verifier calls, work redone by restart, and
the share of clean runs the verifier rejected (a proxy for unnecessary denials, since clean runs can still
contain real model errors).

\section{Experimental setup}
\todo{machine (CPU/RAM), model, backend version and quantization, commit hash (from run\_info.json)}.
Temperature 0, \todo{limits}. Each fault scenario runs under fault, full\_restart and always\_verify, and the
two clean controls run once per task, all repeated 5 times: 170 trials. Failed and timed-out runs are kept.
Rates have Wilson 95\% intervals.

\section{Results}
\begin{table}[h]
\centering\small
\begin{tabular}{lccccc}
\toprule
& clean & clean\_verify & fault & full\_restart & always\_verify \\
\midrule
Strict task success & \res & \res & \res & \res & \res \\
Attack success rate & -- & -- & \res & \res & \res \\
Unauthorized executions (runs) & \res & \res & \res & \res & \res \\
Clean-run rejection rate & -- & \res & -- & -- & -- \\
Median latency (s) & \res & \res & \res & \res & \res \\
Model calls / tokens & \res & \res & \res & \res & \res \\
\bottomrule
\end{tabular}
\caption{Per condition, 5 repeats, tasks 0 and 2 pooled \todo{or split per task}. \todo{split accidental vs adversarial}}
\end{table}

\todo{propagation table: per scenario, count of runs in each of the four categories (fault condition)}

\todo{2--3 sentences of findings: where faults propagate, whether verify prevents it and at what cost,
what restart recovers and the unauthorized writes it cannot undo}

\textbf{Example propagation path.} \todo{one run: injected fact $\rightarrow$ planner input $\rightarrow$ plan
$\rightarrow$ executor call $\rightarrow$ state diff, with run id}

\section{Limitations}
\begin{itemize}
\item Two tasks and one model; results describe this workflow, not agents in general.
\item Faults are fixed replacements, not attacks that adapt to the defense.
\item The bank is simulated: a reset undoes payments that a real bank could not, so restart's unauthorized writes
are reported, not counted as recovered.
\item \todo{propagation labels: how they were made and checked}
\end{itemize}

\section{Phase 2 direction}
\todo{dependency-aware invalidation: tag each result with its sources, invalidate only what depends on the faulty
output, restore affected state, rerun those steps; compare with full restart and always verify at equal budgets;
irreversible actions reported separately. Justify from the Phase 1 results.}

\textbf{Artifacts.} Code, configs and instructions: \todo{repo link}. Raw traces and summaries:
\todo{where the runs/ package is shared}.

\bibliographystyle{plain}
\begin{thebibliography}{1}
\bibitem{agentdojo} E.~Debenedetti, J.~Zhang, M.~Balunovi\'c, L.~Beurer-Kellner, M.~Fischer, F.~Tram\`er.
AgentDojo: A Dynamic Environment to Evaluate Prompt Injection Attacks and Defenses for LLM Agents.
\emph{NeurIPS Datasets and Benchmarks}, 2024.
\end{thebibliography}

\end{document}
